\section{壳上 \texorpdfstring{$\Delta e(u)$}{Delta e(u)} 的计算}
\stfhIndex{I220}{ke^shang}{壳上（on shell）}
\label{sec:calculations-delta-energy-on-shell}

我们将对 $J$-全纯映射的调和能量密度函数
\[
 e(u):=|du|^2
\]
作一些几何计算。这将导出柯西--黎曼方程解的椭圆估计；更细致地研究伪全纯曲线时需要这些估计。

为了最大限度地发挥典范联络的优势，下面的计算采用复化记号。这里涉及殆凯勒流形 $(M,\omega,J)$ 和黎曼曲面 $(\Sigma,j)$。

具体而言，设 $(\Sigma,j)$ 是一个开黎曼曲面，并配备与给定复结构 $j$ 相容的埃尔米特度量 $h$。复化 $T_{\mathbb C}M=TM\otimes\mathbb C$ 具有分解
\begin{equation}
 T_{\mathbb C}\Sigma=T^{(1,0)}\Sigma\oplus T^{(0,1)}\Sigma,
 \qquad
 T_{\mathbb C}M=T^{(1,0)}M\oplus T^{(0,1)}M.
 \label{eq:complexified-tangent-decompositions}
\end{equation}
把导数 $du$ 复线性地延拓为丛映射
\[
 du:T_{\mathbb C}\Sigma\longrightarrow T_{\mathbb C}M
\]
后，它可表示成自然的分块对角矩阵
\begin{equation}
 du=
 \begin{pmatrix}
  (\partial u)'&(\partial u)''\\
  (\overline\partial u)'&(\overline\partial u)''
 \end{pmatrix},
 \label{eq:du-complex-block-matrix}
\end{equation}
其中，$(\cdot)'$ 是其像中 $(\cdot)\in T_{\mathbb C}M$ 的 $(1,0)$ 部分，而 $(\cdot)''$ 是 $(0,1)$ 部分。这里
\begin{equation}
 \partial u\in\Omega^{(1,0)}(\Sigma;u^*T_{\mathbb C}M),
 \qquad
 \overline\partial u\in\Omega^{(0,1)}(\Sigma;u^*T_{\mathbb C}M)
 \label{eq:partial-u-form-types}
\end{equation}
分别是把 $du\in\Omega^1(\Sigma;u^*T_{\mathbb C}M)$ 看作 $\Sigma$ 上取值于 $T_{\mathbb C}M$ 的一形式时，它的 $(1,0)$ 与 $(0,1)$ 分量。这里，我们把切向映射 $du:T_{\mathbb C}\Sigma\to T_{\mathbb C}M$ 看作一形式。

另一方面，回忆\emph{实}算子 $du$ 有分解
\[
 du=\partial_Ju+\overline\partial_Ju.
\]
类似地，把 $\partial_Ju$ 和 $\overline\partial_Ju$ 复线性地延拓为从 $T_{\mathbb C}\Sigma$ 到 $T_{\mathbb C}M$ 的映射。

\begin{lemma}\label{lem:j-partial-complex-block-forms}
在切丛 $T\Sigma$ 的复化 $T_{\mathbb C}M$ 中，作为取值于 $u^*T_{\mathbb C}M$ 的 $\Sigma$ 上一形式，
$\partial_Ju,\overline\partial_Ju:T_{\mathbb C}\Sigma\to T_{\mathbb C}M$ 具有形式
\[
 \partial_Ju=
 \begin{pmatrix}(\partial u)'&0\\0&(\overline\partial u)''\end{pmatrix},
 \qquad
 \overline\partial_Ju=
 \begin{pmatrix}0&(\partial u)''\\(\overline\partial u)'&0\end{pmatrix}.
\]
特别地，$\partial_Ju$ 保持向量的型。
\end{lemma}

\begin{proof}
对所有 $\xi\in T_{\mathbb C}\Sigma$，可以写成
\[
 du(\xi)=
 \left(\frac{du(\xi)-iJ\cdot du(\xi)}2\right)
 +\left(\frac{du(\xi)+iJ\cdot du(\xi)}2\right),
\]
这遵从分解 $T_{\mathbb C}M=T^{(1,0)}M\oplus T^{(0,1)}M$。

现在设 $\xi\in T^{(1,0)}\Sigma$ 并计算
\begin{equation}
 \partial_Ju(\xi)
 =\frac{du(\xi)-J\cdot du\cdot j(\xi)}2
 =\frac{du(\xi)-iJ\cdot du(\xi)}2,
 \label{eq:partial-j-u-one-zero-vector}
\end{equation}
其中使用了 $j\xi=i\xi$。右端恰好是 $du(\xi)$ 到 $T^{(1,0)}M$ 的投影。计算得
\[
 J\left(\frac{du(\xi)-iJ\cdot du(\xi)}2\right)
 =i\left(\frac{du(\xi)-iJ\cdot du(\xi)}2\right).
\]
另一方面，对 $\xi\in T^{(0,1)}\Sigma$，有
\[
 \partial_Ju(\xi)
 =\frac{du(\xi)-J\cdot du\cdot j(\xi)}2
 =\frac{du(\xi)+iJ\cdot du(\xi)}2\in T^{(0,1)}M,
\]
它是 $du(\xi)$ 的 $T^{(0,1)}M$ 投影。因此
\[
 \partial_Ju=
 \begin{pmatrix}(\partial u)'&0\\0&(\overline\partial u)''\end{pmatrix}.
\]
$\overline\partial_Ju$ 的公式同理得到。证明完毕。
\end{proof}

$\Omega^1(u^*T_{\mathbb C}M)$ 分解为直和
\[
 \Omega^1(u^*T_{\mathbb C}M)
 =\Omega^{(0,1)}(u^*T_{\mathbb C}M)
 \oplus\Omega^{(0,1)}(u^*T_{\mathbb C}M),
\]
并且由于我们在黎曼曲面上，有
\[
 \Omega^{(2,0)}(u^*T_{\mathbb C}M)=0
 =\Omega^{(0,2)}(u^*T_{\mathbb C}M),
\]
从而
\[
 \Omega^2(u^*T_{\mathbb C}M)=\Omega^{(1,1)}(u^*T_{\mathbb C}M).
\]

考虑 $u^*TM$ 上仍记为 $\nabla$ 的拉回联络。把它与 $(\Sigma,j,h)$ 的埃尔米特联络结合，便在 $T^*\Sigma\otimes u^*TM$ 上诱导一个联络，我们仍记为 $\nabla$。以 $d^\nabla$ 表示协变导数 $\nabla$ 的斜对称化，它由
\[
 d^\nabla(\alpha)(\xi_1,\xi_2)
 =(\nabla_{\xi_1}\alpha)(\xi_2)
 -(\nabla_{\xi_2}\alpha)(\xi_1)
\]
给出，并定义算子
\[
 d^\nabla:\Omega^1(u^*TM)\longrightarrow\Omega^2(u^*TM).
\]
把 $\alpha=du$ 看作 $\Omega^1(u^*TM)$ 中的一形式，就有下面的引理。

\begin{lemma}\label{lem:covariant-exterior-derivative-du-torsion}
作为取值于 $u^*TM$ 的二形式，有
\begin{equation}
 d^\nabla(du)=u^*T,
 \label{eq:d-nabla-du-pullback-torsion}
\end{equation}
其中 $T$ 是 $\nabla$ 的挠率张量。
\end{lemma}

\begin{proof}
对给定的 $\xi_1,\xi_2\in\Gamma(T\Sigma)$，计算
\[
\begin{split}
 d^\nabla(du)(\xi_1,\xi_2)
 &=\nabla_{\xi_1}(du)(\xi_2)-\nabla_{\xi_2}(du)(\xi_1)\\
 &=\nabla_{du(\xi_1)}(du(\xi_2))-du(\nabla_{\xi_1}\xi_2)
   -\bigl(\nabla_{du(\xi_2)}(du(\xi_1))-du(\nabla_{\xi_2}\xi_1)\bigr)\\
 &=\nabla_{du(\xi_1)}(du(\xi_2))-\nabla_{du(\xi_2)}(du(\xi_1))
   -\bigl(du(\nabla_{\xi_1}\xi_2)-du(\nabla_{\xi_2}\xi_1)\bigr)\\
 &=\nabla_{du(\xi_1)}(du(\xi_2))-\nabla_{du(\xi_2)}(du(\xi_1))-du([\xi_1,\xi_2])\\
 &=\nabla_{du(\xi_1)}(du(\xi_2))-\nabla_{du(\xi_2)}(du(\xi_1))
   -[du(\xi_1),du(\xi_2)]\\
 &=T(du(\xi_1),du(\xi_2))=u^*T(\xi_1,\xi_2),
\end{split}
\]
这就完成了证明。
\end{proof}

把算子 $du$ 与 $T$ 在复化 $T_{\mathbb C}\Sigma$ 和 $T_{\mathbb C}\Sigma$ 上复线性地延拓。

\begin{corollary}\label{cor:j-holomorphic-pullback-torsion-vanishes}
假设 $u$ 是 $J$-全纯的，即 $\overline\partial_Ju=0$。则 $u^*T=0$，从而 $d^\nabla(du)=0$。
\end{corollary}

\begin{proof}
计算
\[
 u^*T(\eta,j\eta)
 =T(du(\eta),du(j\eta))
 =T(du(\eta),Jdu(\eta))=0.
\]
第二个等号使用了 $du\circ j=J\circ du$，第三个等号使用了定理 7.1.5。证明完毕。
\end{proof}

利用推论 7.3.3，得到下述结论。

\begin{theorem}\label{thm:partial-u-harmonic}
设 $u$ 是 $J$-全纯的。则 $\partial u$ 是调和的，即满足
$\Delta(\partial u)=0$。
\end{theorem}

\begin{proof}
设 $*$ 是 $\Sigma$ 上的霍奇星算子。对
$\xi\in(T^*\Sigma)^{(1,0)}$，有 $*\xi=-i\xi$。因此
\begin{equation}
 *\partial u=-i\partial u.
 \label{eq:hodge-star-partial-u}
\end{equation}
利用伴随算子的公式
\[
 (d^\nabla)^*=-*d^\nabla*
\]
并把它应用于 $\partial u$，得到
\[
 (d^\nabla)^*(\partial u)
 =-*d^\nabla*(\partial u)=i*d^\nabla(\partial u).
\]
由推论 7.3.3，这等于零，因为
$d^\nabla\partial u=d^\nabla du|_{T^{(1,0)}\Sigma}$。

现在考虑霍奇拉普拉斯算子
\[
 \Delta=\Delta^u:=d^\nabla(d^\nabla)^*+(d^\nabla)^*d^\nabla
\]
并计算
\begin{equation}
 \Delta(\partial u)
 =d^\nabla(d^\nabla)^*(\partial u)
 +(d^\nabla)^*d^\nabla(\partial u)=0.
 \label{eq:partial-u-hodge-laplacian-zero}
\end{equation}
证明完毕。
\end{proof}

现在回忆博赫纳--魏岑伯克公式
\stfhIndex{I064}{bo^he^na^wei^cen^bo^ke^gong^shi}{博赫纳--魏岑伯克公式（Bochner–Weitzenböck formula）}
\begin{equation}
 \Delta(\partial u)
 =\nabla^*\nabla(\partial u)+K(\partial u)+R(du,du)\partial u,
 \label{eq:bochner-weitzenbock-partial-u}
\end{equation}
其中 $\nabla^*\nabla$ 是迹拉普拉斯算子，$K$ 是 $(\Sigma,h)$ 的高斯曲率，而 $R$ 是 $M$ 上联络 $\nabla$ 的截面曲率，并复线性地延拓到 $T_{\mathbb C}M$。（关于一般黎曼流形上向量值微分形式的魏岑伯克公式之完整推导，见附录 A.1。）结合定理 7.3.4 与\eqref{eq:bochner-weitzenbock-partial-u}，已经得到
\begin{equation}
 \nabla^*\nabla(\partial u)
 =\Delta(\partial u)-K\partial u-R(du,du)\partial u
 =-K\partial u-R(du,du)\partial u.
 \label{eq:trace-laplacian-partial-u}
\end{equation}

对一个 $J$-全纯映射 $u:(\Sigma,j)\to(M,J)$，计算能量密度 $e(u)$ 的拉普拉斯算子，其中 $\Sigma$ 可以具有非空边界 $\partial\Sigma$。回忆，对 $J$-全纯映射有恒等式
\[
 e(u)=|du|^2=|\partial u|^2+|\overline\partial u|^2=2|\partial u|^2.
\]
\emph{换言之，能量密度可以完全用取值于 $T^{(1,0)}M$ 的 $(1,0)$-形式 $\partial u$ 表示。}因此，利用这个恒等式，我们计算 $\Delta|\partial u|^2$ 而非 $\Delta|du|^2$。注意到 $du$ 是实算子，可导出
\begin{equation}
\begin{split}
 \frac12\Delta e(u)
 &=\Delta\langle\partial u,\partial u\rangle\\
 &=\langle\nabla^*\nabla\partial u,\partial u\rangle
 -2|\nabla\partial u|^2
 +\langle\partial u,\nabla^*\nabla\partial u\rangle\\
 &=-2K\langle\partial u,\partial u\rangle
 -2\langle R(du,du)\partial u,\partial u\rangle
 -2|\nabla\partial u|^2.
\end{split}
 \label{eq:laplacian-energy-density}
\end{equation}
所以
\begin{equation}
 |\nabla du|^2=2|\nabla\partial u|^2
 =-\frac12\Delta e(u)
 -2K\langle\partial u,\partial u\rangle
 -2\langle R(du,du)\partial u,\partial u\rangle.
 \label{eq:nabla-du-energy-density}
\end{equation}
这里，第一个等号应用了下面的引理。

\begin{lemma}\label{lem:nabla-du-partial-u-norm}
有 $2|\nabla\partial u|^2=|\nabla du|^2$。
\end{lemma}

\begin{proof}
写成 $\nabla du=\nabla(\partial u+\overline\partial u)$。设
$\{e_1,e_2\}$ 是 $T\Sigma$ 的一个正交归一标架。则
\[
 |\nabla(\partial u)|^2
 =|\nabla_{e_1}(\partial u)|^2+|\nabla_{e_2}(\partial u)|^2.
\]
另一方面，由定义 $\partial u=du|_{T^{(1,0)}\Sigma}$，有
\[
 |\nabla_{e_1}(\partial u)|^2
 =|\nabla_{e_1}(du)(\xi)|^2,
 \qquad
 \xi=\frac1{\sqrt2}(e_1-ie_2).
\]
类似地，有
\[
 |\nabla_{e_1}(\overline\partial u)|^2
 =|\nabla_{e_1}(du)(\overline\xi)|^2.
\]
但是
\[
 \nabla_{e_1}(du)(\overline\xi)
 =\nabla_{e_1}(du(\overline\xi))-du(\nabla_{e_1}(\overline\xi))
 =\overline{\nabla_{e_1}(du)(\xi)},
\]
因而 $|\nabla_{e_1}(\partial u)|^2=|\nabla_{e_1}(\overline\partial u)|^2$。类似地，也有
$|\nabla_{e_2}(\partial u)|^2=|\nabla_{e_2}(\overline\partial u)|^2$。这证明
$|\nabla(\partial u)|^2=|\nabla\overline\partial u|^2$。回忆
$|\nabla du|^2=|\nabla\partial u|^2+|\nabla\overline\partial u|^2$，证明完毕。
\end{proof}

当 $\partial\Sigma=\varnothing$ 时，在 $\Sigma$ 上积分\eqref{eq:nabla-du-energy-density}，得到
\begin{equation}
 \int_\Sigma|\nabla du|^2
 =-2\int_\Sigma\langle K\partial u,\partial u\rangle
 -2\int_\Sigma\langle R(du,du)\partial u,\partial u\rangle.
 \label{eq:closed-surface-integrated-bochner}
\end{equation}

\begin{proposition}\label{prop:closed-j-holomorphic-w12-bound}
对任意满足 $\partial\Sigma=\varnothing$ 的 $J$-全纯映射 $u:\Sigma\to M$，都有不等式
\begin{equation}
 \int_\Sigma|\nabla du|^2
 \leq\int_\Sigma\max_{z\in\Sigma}(-K)(z)|du|^2
 +\frac12\int_\Sigma\max_{x\in M}|R|(x)|du|^4.
 \label{eq:closed-j-holomorphic-integral-bound}
\end{equation}
\end{proposition}

\begin{proof}
这由\eqref{eq:boundary-integrated-bochner}和恒等式
$|du|^2=2|\partial u|^2$ 得到。
\end{proof}
